%% CPP 2027 Submission — 0pat Exercise XXVIII: Observing the Critical Line in Split-Complex Coordinates
\documentclass[sigplan,10pt,anonymous,review]{acmart}
\settopmatter{printfolios=true,printacmref=false}

\AtBeginDocument{%
  \providecommand\BibTeX{{\normalfont B\scshape ib\TeX}}}

\settopmatter{printacmref=false}
\renewcommand\footnotetextcopyrightpermission[1]{}
\pagestyle{plain}

\usepackage{microtype}
\usepackage{booktabs}
\usepackage{amsmath}
% XeTeX (tectonic): acmart loads newtxmath/libertine which already defines \Bbbk;
% reset it so amssymb's definition is accepted.
\let\Bbbk\relax
\usepackage{amssymb}

\begin{document}

\title{Observing the Critical Line in Split-Complex Coordinates\\
(Exercise XXVIII)}

\author{Anonymous}
\affiliation{%
  \institution{Anonymous}
}
\email{anonymous@example.org}

\begin{abstract}
We observe the critical line of the Riemann direction through the
split-complex frame (Exercise XXVII): the imaginary-radius orthogonal
Euler circles. The observation yields three machine-checked results.
First, the pseudo-norm of a critical-line point is
$\|1/2 + tj\|^2 = 1/4 - t^2$: the real-part condition
$\mathrm{Re}(s) = 1/2$ becomes, in split coordinates, an imaginary-part
threshold $|t| = 1/2$ --- the two axes exchange roles. Second, the
threshold is structural: points with $|t| < 1/2$ are timelike (inside
the radius-1 circle) and points with $|t| > 1/2$ are spacelike (inside
the radius-$i$ circle); since the known zeros have $|t| \ge 14.13 >
1/2$, every known critical-line zero is spacelike in the split frame.
Third, the reciprocal map sends the critical line to the hyperbola
$(a-1)^2 - b^2 = 1$ --- the split dual of the critical circle
$(a-1)^2 + b^2 = 1$ already proved in the complex frame (C019--C022):
the same map, two hosts, circle $\leftrightarrow$ hyperbola. Six
theorems, zero errors, zero \texttt{sorry}, mathlib v4.32.2.
\end{abstract}

\keywords{Lean, formalization, split-complex, critical line, dual structure, 0pat}

\maketitle

\section{The observation}

The split-complex frame (Exercise XXVII) hosts the imaginary-radius
orthogonal Euler circles: $a + bj$, $j^2 = +1$, pseudo-norm
$\|a + bj\|^2 = a^2 - b^2$. The radius-1 circle is the hyperbola
$a^2 - b^2 = 1$; the radius-$i$ circle is the conjugate hyperbola
$a^2 - b^2 = -1$. We now observe the critical line of the Riemann
direction --- the line $\mathrm{Re}(s) = 1/2$ --- through this frame:
its points are $\frac12 + tj$, $t \in \mathbb{R}$.

The observation is three theorems.

\paragraph{The pseudo-norm of the critical line.}
$\|\frac12 + tj\|^2 = \frac14 - t^2$. The real-part condition
$\mathrm{Re}(s) = 1/2$, which defines the critical line in the complex
frame, becomes in the split frame an imaginary-part threshold:
$|t| = 1/2$. The two axes exchange roles --- the split metric swaps the
meaning of the coordinates.

\paragraph{The threshold.} Points with $t^2 < 1/4$ are timelike
($\|z\|^2 > 0$, inside the radius-1 circle); points with $t > 1/2$ are
spacelike ($\|z\|^2 < 0$, inside the radius-$i$ circle). Since the known
zeros of $\zeta$ on the critical line have imaginary parts
$|t| \ge 14.13 > 1/2$, every known critical-line zero is spacelike in
the split frame --- the radius-$i$ region is where the zeros live. The
threshold $\frac12$ is not chosen; it is the value of the frame.

\paragraph{The dual of the critical circle.} The reciprocal map
$z \mapsto z^{-1}$ (the split inverse, valid when $\|z\|^2 \ne 0$) sends
the critical line to the hyperbola $(a-1)^2 - b^2 = 1$. In the complex
frame, the same reciprocal map sends the critical line to the circle
$(a-1)^2 + b^2 = 1$ --- the critical circle, proved in the Riemann chain
(C019--C022). The split observation is the exact dual: the same map,
two hosts, circle $\leftrightarrow$ hyperbola. The center $(1, 0)$ and
the radius $1$ are unchanged; only the metric changes.

\section{Honest boundary and position}

The content is classical mathematics (split-complex algebra, elementary
computations); the value is observational: the critical-line geometry of
the Riemann direction has a dual structure in the split frame, and the
dual places the zeros (if on the critical line) in the spacelike region
with a structural threshold $\frac12$. No claim is made about the zeros
themselves: the observation concerns the geometry of the line, not the
analysis of $\zeta$. The frame observes; the analysis remains open.

\section{Formalization}

The proof is a single Lean file (\texttt{proof.lean}, 127 lines),
importing only \texttt{Mathlib.Data.Real.Basic} and basic tactics. The
split-complex structure (30 lines, as in Exercise XXVII) is
self-contained. Six theorems, compiled with Lean 4.32.2 / mathlib
v4.32.2, zero errors, zero \texttt{sorry}.

\begin{center}
\begin{tabular}{lll}
\toprule
Theorem & Statement & Proof core \\
\midrule
\texttt{inv\_mul\_self} & $z^{-1} z = 1$ ($\|z\|^2 \ne 0$) & field\_simp; ring \\
\texttt{critical\_line\_pseudonorm} & $\|\frac12 + tj\|^2 = \frac14 - t^2$ & simp; ring \\
\texttt{threshold\_timelike} & $t^2 < 1/4 \Rightarrow \|z\|^2 > 0$ & nlinarith \\
\texttt{threshold\_spacelike} & $t > 1/2 \Rightarrow \|z\|^2 < 0$ & nlinarith \\
\texttt{zero\_observation\_spacelike} & $|t| > 1/2 \Rightarrow \|z\|^2 < 0$ & sign split; nlinarith \\
\texttt{recip\_critical\_line\_hyperbola} & $\mathrm{recip}(\frac12 + tj)$ on $(a-1)^2 - b^2 = 1$ & field\_simp; ring \\
\bottomrule
\end{tabular}
\end{center}

\section{Conclusion}

The critical line, observed through the imaginary-radius Euler circles,
reveals its dual: the circle of the complex frame becomes the hyperbola
of the split frame under the same reciprocal map, and the threshold
$\frac12$ appears as a structural value of the metric. The zeros of
$\zeta$ (if on the critical line) live in the spacelike region, inside
the radius-$i$ circle. The observation is complete; the analysis that
would place the zeros there remains the open span of the bridge.

\bibliographystyle{ACM-Reference-Format}
\bibliography{refs}

\end{document}
